\section{Browsing}
\label{sec:browsing}

\subsection{Names and case}
\label{sec:names}

\begin{req}{NAM-1}
Wherever this specification says that two names are ``the same name'', it means
what the file system means.  On the macOS file system in its default,
case-insensitive form, two names are the same name when they are equal after
Unicode normalization and \emph{full} case folding.  Under full folding,
\code{ss} and the letter sharp s, \code{s} and the long s, \code{k} and the Kelvin
sign, and the ligatures \code{ff}, \code{fi}, \code{fl}, \code{ffi}, \code{ffl}
and \code{st} are the same as their spelled-out forms.  A name that the file system
resolves to a file is that file, and every rule, root check and search treats it as
such.
\end{req}

\begin{req}{NAM-2}
A name that contains characters that would disguise it, namely control characters,
bidirectional overrides and isolates, line and paragraph separators, the byte order
mark, and tag characters, is displayed with each such character replaced by a
visible marker naming its code point (for example \code{\textlangle U+202E\textrangle}).
Ordinary text, including text in right-to-left scripts, is displayed unchanged.
Links and downloads use the real name.
\end{req}

\begin{req}{NAM-3}
A name that is not valid UTF-8, which Linux allows, is listed, found by search,
described and opened like any other.  Each byte of it that is not part of valid
UTF-8 is displayed as a visible marker naming its value (for example
\code{caf\textlangle 0xE9\textrangle.txt}), so two names that differ only in such
bytes look different.  ``Copy path'' gives such a path quoted for the shell, with the
byte escaped (\code{\$'caf\textbackslash xE9.txt'}), which bash and zsh read back as
the real name.  The interface carries the name as \rid{API-10} describes.
\end{req}

\subsection{Roots and reachability}

\begin{req}{ROOT-1}
Only the roots (\rid{CLI-2}, \rid{CLI-3}) are reachable.  A request for anything
outside them, including through a symbolic link that leads outside, is answered as a
path that does not exist.  Whether a path is inside a root is decided by the
identity of the folder, not by how the path is spelled.  (On a case-sensitive
volume, a root named \code{Public} does not admit a different folder named
\code{public}.)
\end{req}

\begin{req}{ROOT-2}
When there is more than one root, the interface offers a way to switch between
them.  The box \emph{Go to path} (\code{Alt+G}) accepts an absolute path, or one
beginning \code{\~{}/}, and resolves \code{.}, \code{..}, doubled and trailing
slashes.  A path outside the roots produces a message that says the folder is not
served and lists the roots.  \code{\~{}name} (another user's home) is refused.
\end{req}

\subsection{Listing}

\begin{req}{LST-1}
Opening a folder lists its entries with the columns \emph{Name}, \emph{Size},
\emph{Modified}, \emph{Kind} and \emph{Permissions}, and an optional last column,
\emph{Attributes} (\rid{META-3}).  A folder has no size shown.  Size is right-aligned.
\end{req}

\begin{req}{LST-2}
The listing is delivered incrementally, so the first rows appear before a very
large folder has been read, and only the rows on screen are drawn.  A folder of
100{,}000 entries remains responsive (Section~\ref{sec:limits}).
\end{req}

\begin{req}{LST-3}
Entries are sorted as the tool \code{eza} sorts by default: names compare without
regard to case, digits inside names compare by value (\code{file2} before
\code{file10}), and a leading dot sorts before letters, so \code{.bashrc} comes
first.  Folders are \emph{not} grouped by default.  Clicking a column header sorts
by that column, and clicking it again reverses the order.  Ties fall back to the
name and, if two names still compare equal (\code{file1} and \code{file01}, whose
digits compare by value), to their raw characters, so the order is total and does
not depend on the order in which the file system returned the entries.
\end{req}

\begin{req}{LST-4}
The option \emph{Folders first}, off by default, lists all folders before all files
under every sort column and in both directions.  A folder, having no size, sorts as
smaller than any file when sorting by size.  A modification time that cannot be read
sorts as the oldest.
\end{req}

\begin{req}{LST-5}
A symbolic link is shown with a link marker.  Its size, date and kind are those of
its target.  A link whose target does not exist is shown as broken.  A link whose
target is denied, or lies outside the roots, is not shown at all.
\end{req}

\begin{req}{LST-6}
Sockets, devices and pipes are listed but never opened.
\end{req}

\begin{req}{LST-7}
Entries whose names begin with a dot are hidden until the option \emph{Hidden files}
is turned on.  This option is independent of the ignore file (\rid{RUL-2}) and never
reveals an entry that a rule hides or denies.
\end{req}

\begin{req}{LST-8}
The box \emph{Filter this folder} (key \code{/}) narrows the listing as the user
types, to entries whose names contain the text.  Matching ignores case and
normalization as in \rid{NAM-1}, unless the option \emph{Match case} is on, in which
case only normalization is ignored.  The filter finds the same names that search
finds (\rid{SRC-3}).
\end{req}

\begin{req}{LST-9}
The interface shows the current path as breadcrumbs, each a link to that folder,
and a button that copies the path (also \code{Alt+C}, which copies the selected
entry's path, or the folder's when nothing is selected).
\end{req}

\begin{req}{LST-10}
The location in the browser's address bar identifies the folder (in the fragment,
\code{\#/Users/me/docs}) and optionally an entry to select
(\code{?select=name}), so a folder can be bookmarked and the browser's Back button
works while \texttt{fsb} is running.
\end{req}

\begin{req}{LST-11}
Columns can be resized by dragging their edge (double-click, or Enter on the edge,
fits the column to the visible rows; Left and Right resize by 10 pixels, with Shift
by 50), reordered by dragging a header onto another or with Alt and the arrow keys,
and shown or hidden from the \emph{Columns} menu (Name is always shown).  Widths are
between 60 and 900 pixels.  \emph{Reset columns} restores the defaults.  The menu
closes on an outside click and on Escape.  On a narrow window the table scrolls
sideways and does not hide columns.
\end{req}

\begin{req}{LST-12}
Denied entries never appear in a listing (\rid{RUL-1}), whatever options are set.
The interface shows a badge saying that it is read-only.
\end{req}

\begin{req}{LST-13}
Opening a folder that is itself stored only in the cloud (not downloaded) reports that
plainly and reads nothing, rather than listing it as empty or answering with an error
that does not say why (\rid{SEC-13}).
\end{req}

\subsection{Search}

\begin{req}{SRC-1}
The box \emph{Search subfolders by name} (\code{Alt+S}; Enter runs the search)
finds entries under the current folder whose names contain the text, at any depth.
\end{req}

\begin{req}{SRC-2}
Results stream in as they are found, shallowest first, so the nearest matches
appear immediately.  The search stops at 1{,}000 results or after 500{,}000 entries
have been examined, and the interface says when a limit stopped it.
\end{req}

\begin{req}{SRC-3}
Search compares names as in \rid{NAM-1} (case, accents and normalization are
ignored), unless \emph{Match case} is on.  A search for \code{strasse} finds a file
named with the sharp s.
\end{req}

\begin{req}{SRC-4}
Search never enters or reports a denied, hidden or cloud-only folder, never reports a
denied or hidden entry, and does not follow symbolic links to folders (so it always
terminates).  A cloud-only folder is skipped rather than failing the whole search.
Anything a search reports is also shown when the folder that contains it is listed.
A folder that was queued for searching and, by the time its turn came, had been
replaced by a symbolic link (or lies behind one) is not entered, because its entries
would be judged under a name that is not where they are; it counts as not searched
(\rid{SRC-6}).
\end{req}

\begin{req}{SRC-5}
Choosing a result opens its folder with that entry selected.  While results are
shown, the last breadcrumb, Escape, and a \emph{Clear search} control all return to
the plain listing.
\end{req}

\begin{req}{SRC-6}
A search that could not look at everything it was allowed to look at says so, and its
results are not presented as complete.  This covers a folder that is there and
allowed but could not be opened (a permission error, an input/output error), a
folder that failed partway through being read, and a folder replaced by a symbolic
link as in \rid{SRC-4}.  A denied folder, which looks exactly like a missing one,
and a folder that vanished are not counted, and neither is a cloud-only folder, which
is skipped on purpose (\rid{SRC-4}): completeness is relative to the searchable,
allowed content.
\end{req}
